\documentclass[10pt]{article}
\usepackage{icmlko}

\icmlmeta{IP 파운데이션 월드모델}{974}{삼중쌍이라 부르기 전에 정밀도를 재라}{앞 노트가 「(s,a,o) 삼중쌍 35,641」이라 부른 것의 개체 정밀도는 0.754310 이고, 그 자료를 넣은 유보 예측의 leave-one-source-out $\Delta$ 는 0.002192 (순열 $p$ = 0.465134)다}

\begin{document}

\twocolumn[
\icmlheader

\begin{abstract}
앞 노트는 113GB 한국어 웹 코퍼스를 처음 끝까지 정제해 \textbf{「(s,a,o) 삼중쌍 35,641 행」}을
냈다고 적었다. 🔴 \textbf{그 이름은 아직 못 쓴다} --- 언급이 참인지 \textbf{한 번도 안 쟀기}
때문이다. 이 노트는 무작위 표본 \textbf{696 행}(층 = 제목 토큰 단일/다중 $\times$ 개체 언어
ko/en)을 \textbf{문서 맥락과 함께 읽어} 라벨하고, \textbf{개체 정밀도 0.754310}
(95\% Clopper--Pearson 0.720571 -- 0.785877)와 층 가중 \textbf{유효 삼중쌍 27,499.4}
(분모 35,641 · 77.16\%)를 낸다.
🔴 층별로는 \textbf{언어가 토큰 수보다 크게 가른다} --- 다중·ko 0.929032 $>$ 단일·ko 0.792254
$>$ 다중·en 0.640777 $>$ 단일·en 0.584416 이라, \textbf{「단일 토큰 일치는 곧 잡음」이라는 읽기는
반만 맞다}(단일·ko 는 22,640 행으로 가장 큰 층인데 다중·en 보다 정밀하다).
그 다음에야 \textbf{이 축의 자}를 처음 잰다: 유보를 시간으로 자르면 새 원천은 \textbf{학습에만}
들어가고 \textbf{유보 집합이 팔마다 글자 그대로 같다}. 그 설계에서
\textbf{leave-one-source-out $\Delta\rho$ = 0.002192} (순열 $p$ = 0.465134 · 귀무 95 분위
0.034364) --- \textbf{채택 실패}이고, 정밀도 게이트를 건 팔은 \textbf{-0.044382} 로 더 나쁘다.
🔴 \textbf{그러나 같은 주행이 다른 축 하나를 움직인다} --- \textbf{도메인별 유보 성능의 최저값}이
\textbf{-0.345988 $\to$ -0.235423 $\to$ -0.158134} 로 단조로 오르는데 \textbf{평균은 거의 안 움직인다}
(0.215083 $\to$ 0.211554). \textbf{행을 늘리는 것은 전체를 올리는 일이 아니라 바닥을 올리는 일이었다.}
\end{abstract}
]

\section{이름이 자를 앞서면 안 된다}

앞 노트는 \textbf{없던 것을 만들었다} --- 그 코퍼스에서 온 행이 0 이었고 35,641 이 되었다.
그것은 참이다. 🔴 \textbf{그러나 그 행을 「삼중쌍」이라 부른 순간 두 가지가 같이 참이 되어야 한다}:
상태 $s$ 와 결과 $o$ 가 그 개체의 것이고, 액션 $a$ 가 \textbf{그 개체를 실제로 가리키는} 것이다.
앞의 것은 정의로 보장된다(둘 다 같은 위키 계열을 자른 것이다). 뒤의 것은 \textbf{측정이 필요하다.}

이 노트의 첫 결정은 \textbf{순서}다. 정밀도를 모르는 행으로 기여도를 재면,
$\Delta$ 가 0 근처로 나와도 \textbf{자료가 안 도운 것인지 자료가 틀린 것인지 원리상 못 가른다.}
그래서 사전등록에 \textbf{「정밀도를 안 내면 실패」}와 \textbf{「정밀도 산출물이 없으면 학습 러너가
터진다」}를 같이 박았다.

\claimx{규모를 보고한 뒤 정밀도를 보고하는 순서는 결과의 해석을 바꾼다}{정밀도를 재기 전과 후에 $\Delta$ 의 해석이 같다면 이 주장은 떨어진다. 여기서는 77.16\% 라는 수가 있어야 「자료가 쓰레기라서」라는 설명을 뺄 수 있다.}

\section{표본과 라벨}

\textbf{층}은 행의 필드에서만 정한다 --- 맞은 제목의 공백 토큰 수(단일/다중)와 개체 언어(ko/en).
층마다 무작위로 뽑고(씨앗 974), 곁 표본으로 층 없는 단순무작위 300 행을 따로 뽑아 대조했다.
표본 행의 문서 id 로 \textbf{shard 를 다시 흘려 읽어} 본문을 되찾고(657 문서 · 못 찾은 문서
0), 맞은 자리 앞뒤를 잘라 산출물에 남겼다 --- \textbf{라벨을 재검할 수 있어야} 하기 때문이다.

라벨은 네 값이다. \textbf{참}(그 자리가 그 개체를 가리킨다) · \textbf{거짓·일반어}(그 위키 문서
제목의 \emph{일반어 뜻}을 가리키고 개체는 안 가리킨다) · \textbf{거짓}(제3의 것) · \textbf{모름}.
🔴 셋째 값이 필요한 이유는 실측에서 나왔다 --- 웹툰 개체 \texttt{블루투스} 의 상태 계열은
\textbf{무선 규격} 문서의 조회수다. 그런 행에서 「블루투스 4.0」은 \textbf{문서에는 맞고 개체에는
틀리다}. 그래서 두 정밀도를 같이 적는다: \textbf{개체 정밀도 0.754310} 와
\textbf{문서 정밀도 0.831897}.

\claimx{개체 기준 정밀도와 문서 기준 정밀도를 나누지 않으면 언급 자료의 오류를 잘못 귀속한다}{두 값이 같으면 이 주장은 떨어진다. 여기서는 0.754310 대 0.831897 로 54 행이 그 사이에 있다.}

\section{정밀도}

주 자는 \textbf{모름을 실패로 세는} 개체 정밀도다: \textbf{525 / 696 = 0.754310}
(95\% CP 0.720571 -- 0.785877). 곁 표본(단순무작위 300)은 0.786667 로 구간이 겹친다.
층 가중으로 되돌린 \textbf{유효 삼중쌍은 27,499.4}(붓스트랩 95\% 26,228.8 -- 28,677.4).

층별 정밀도는 \textbf{단일·en 0.584416} · \textbf{다중·en 0.640777} · \textbf{단일·ko 0.792254}
· \textbf{다중·ko 0.929032} 다. 사전등록은 「가장 낮은 층이 단일·en 일 것」을 맞혔고
「정밀도가 0.75 미만일 것」과 「유효 삼중쌍이 25,000 미만일 것」을 \textbf{틀렸다}.
🔴 \textbf{두 반증 모두 앞 노트에게 유리한 쪽이다. 그대로 적는다.}

\claimx{토큰 수보다 언어가 언급 정밀도를 크게 가른다}{단일·ko 가 다중·en 보다 낮으면 이 주장은 떨어진다. 여기서는 0.792254 대 0.640777 다.}

\section{게이트 --- 그리고 성립하지 않은 게이트 하나}

네 게이트를 걸고 \textbf{라벨로 채점}했다: 단일 토큰 라틴 제목 · 토큰화가 원문을 뭉개는 제목 ·
위키 자기 순환 host · 레포트/덤프 사이트. 넷을 다 걸면 \textbf{11,332 / 35,641 (31.79\%)}
가 떨어지고, 떼는 쪽 정밀도 \textbf{0.678967} 가 남는 쪽 \textbf{0.802353} 보다 낮다.

🔴 \textbf{그런데 위키 자기 순환은 게이트로 성립하지 않는다} --- 무는 쪽 정밀도가
\textbf{0.804348} 로 남는 쪽 \textbf{0.750769} 보다 \emph{높다}. 위키 미러 문서는 맞는 제목을
그대로 나열하니 \textbf{언급이 오히려 정확하다}. \textbf{자기 순환은 정밀도 문제가 아니라 독립성
문제이고, 정밀도 자로는 원리상 못 잡는다.}

\claimx{자기 순환 원천은 정밀도 자로 걸러지지 않는다}{위키 host 를 무는 행의 정밀도가 남는 행보다 낮으면 이 주장은 떨어진다. 여기서는 0.804348 대 0.750769 로 반대다.}

\section{문턱의 손실을 같은 표본에서 센다}

한글 제목 길이 문턱(4자)은 진짜 개체도 뗀다. 문턱을 3 으로 내리면 제목 133 개가 새로
들어오고 같은 표본 문서에서 짝이 \textbf{1,194 개} 는다. 🔴 \textbf{그중 840 개가
`원피스' 하나}이고 뒤이어 `네이버' · `고양이' · `거짓말' 같은 \textbf{세 글자 보통명사}가 온다.
문턱이 뗀 진짜 개체(`원펀맨')는 \textbf{12 짝}이다. \textbf{문턱은 손실보다 이득이 크다.}

\section{이 축의 자 --- leave-one-source-out}

\textbf{설계}. 유보를 \textbf{시간}으로 자른다. 새 원천의 문서는 유보 창보다 앞선 것뿐이라
\textbf{유보에 0 행} 들어간다 --- 그래서 \textbf{유보 집합이 세 팔에서 글자 그대로 같고
다른 것은 학습뿐}이다. 그것이 「원천 하나를 빼면 유보 예측이 얼마나 달라지나」의 깨끗한 꼴이다.
특징 · 능형 · 순위상관 · 묶음 · 순열 귀무는 \textbf{앞 노트들의 생산 함수를 그대로 부른다}.

\textbf{결과}. 학습 행 1,581 $\to$ 37,222 로 \textbf{35,641 행}을 더해도
\textbf{묶음 $\Delta\rho$ = 0.002192} · 순열 $p$ = \textbf{0.465134} · 양수 도메인 4 / 7 다.
귀무 95 분위가 0.034364 이니 \textbf{관측은 이 설계가 가를 수 있는 크기 밑}이다.
정밀도 게이트를 건 팔은 \textbf{-0.044382}($p$ = 0.954511)로 \textbf{더 나쁘다}.

🔴 정직한 문장은 \textbf{「좋아진다도 나빠진다도 이 분모로는 못 말한다」}이다.
🔴 \textbf{그러나 앞 노트와 달라진 것이 하나 있다} --- 정밀도를 쟀기 때문에
\textbf{「자료가 쓰레기라서」는 이제 답이 아니다}. 행의 77.16\% 가 진짜 언급이다.

\claimx{언급 정밀도가 높아도 규모만으로는 유보 예측이 안 오른다}{$\Delta$ 가 귀무 95 분위를 넘으면 이 주장은 떨어진다. 여기서는 0.002192 대 0.034364 다.}

\section{바닥이 오른다 --- 평균은 안 오른다}

같은 주행이 \textbf{도메인별 유보 성능}을 낸다. \textbf{최저값}은
\textbf{-0.345988 $\to$ -0.235423 $\to$ -0.158134} 로 \textbf{단조로 오르고}, 최저 도메인은 세 팔 모두
같다. 🔴 \textbf{그런데 평균은 0.215083 $\to$ 0.213073 $\to$ 0.211554 로 거의 안 움직인다.}

\textbf{묶음 $\Delta$ 는 0 근처인데 바닥은 오른다} --- 이 둘은 모순이 아니다. 도메인별 $\Delta$ 는
\textbf{Δ$^+$ 0.110565 와 최대 Δ$^-$ -0.206897} 로 갈리고, 유보 행 가중 합이 서로를 지운다.
\textbf{규모를 더하는 것은 평균을 올리는 일이 아니라 재분배}였고, \textbf{그 재분배가 바닥에 유리하다.}

\claimx{데이터 규모를 늘리는 이득은 평균이 아니라 최저 도메인에서 나타난다}{최저값이 안 오르거나 평균이 같이 오르면 이 주장은 떨어진다. 여기서는 최저 -0.345988 $\to$ -0.158134 이고 평균 0.215083 $\to$ 0.211554 다.}

\section{한계}

라벨은 \textbf{한 사람}이 붙였고 \textbf{둘째 채점자가 없다}. 맞은 자리 앞뒤 맥락만 보았고 문서 전문은
안 읽었다. 새 원천의 문서는 유보 창보다 \textbf{여러 해 앞선} 것이라 \textbf{분포 이동}이 팔 사이에
섞여 있고, 이 설계로는 「원천이 나빠서」와 「시대가 달라서」를 못 가른다. 그리고 읽은 shard 는 앞
노트와 같은 여덟 개다 --- \textbf{전량이 아니다}.

\end{document}
